% ============================================================
% AAAI-2027 Submission
% "Hierarchical Decomposition of Context: A Three-Level
%  Multi-Agent Architecture for Autonomous Software Engineering"
% ============================================================
\documentclass[letterpaper]{aaai25}  % Use aaai25 style (AAAI 2025/2026/2027 compatible)

% ---- Required AAAI packages --------------------------------
\usepackage{times}
\usepackage{helvet}
\usepackage{courier}
\usepackage[hyphens]{url}
\usepackage{graphicx}
\usepackage{amsmath,amssymb,amsthm}
\usepackage{algorithm}
\usepackage{algpseudocode}
\usepackage{booktabs}
\usepackage{multirow}
\usepackage{xspace}
\usepackage{microtype}
\usepackage[hidelinks]{hyperref}
\usepackage{natbib}

% ---- Theorem environments ----------------------------------
\newtheorem{theorem}{Theorem}
\newtheorem{proposition}{Proposition}
\newtheorem{definition}{Definition}
\newtheorem{corollary}{Corollary}
\newtheorem{lemma}{Lemma}

% ---- Macros ------------------------------------------------
\newcommand{\sysname}{\textsc{SubAgent Manager}\xspace}
\newcommand{\Hm}{\mathcal{H}}      % High Manager (Level 1)
\newcommand{\Dm}{\mathcal{D}}      % Domain Manager (Level 2)
\newcommand{\Wa}{\mathcal{W}}      % Worker Agent (Level 3)
\newcommand{\ctx}{\mathbf{c}}      % context vector
\newcommand{\tools}{\mathcal{T}}   % tool set
\newcommand{\plan}{\Pi}            % plan
\newcommand{\patch}{\delta}        % patch
\newcommand{\verify}{\mathcal{V}}  % verification predicate

% ---- Paper metadata ----------------------------------------
\title{Hierarchical Decomposition of Context:\\
A Three-Level Multi-Agent Architecture for Autonomous Software Engineering}

\author{
  % [Blind review submission — author information omitted]
  Anonymous Authors\\
  Anonymous Institution\\
  \texttt{anonymous@anonymous.edu}
}

\begin{document}

\maketitle

% ============================================================
\begin{abstract}
% ============================================================
We present \sysname, a three-level hierarchical multi-agent framework for
autonomous software engineering, and evaluate it on SWE-bench Lite
\cite{jimenez2024swebench}---the canonical benchmark of 300 real-world
GitHub issue resolution tasks drawn from 12 open-source Python repositories.
Our architecture decomposes the software debugging workflow into three
hierarchical tiers: (1) a singleton \emph{High Manager} ($\Hm$) that
enforces a four-phase verification pipeline (Analysis $\to$ Reproduction
$\to$ Patch $\to$ Verification), (2) three \emph{Domain Managers}
($\Dm_{\text{analysis}},\, \Dm_{\text{testing}},\, \Dm_{\text{patch}}$)
that independently plan and delegate domain-specific subtasks, and (3) six
\emph{Worker Agents} ($\Wa_i$) each equipped with a minimal, surgically
scoped tool set.

All inference is performed locally on a single consumer GPU (NVIDIA RTX
5080, 16\,GB GDDR7) using
\texttt{Qwen2.5-Coder-7B-Instruct-AWQ}~\cite{qwen25coder2024} served via a
custom \texttt{vLLM}~\cite{kwon2023vllm} asynchronous engine with 4-bit
Activation-Aware Weight Quantization (AWQ)~\cite{lin2024awq}. On SWE-bench
Lite (all 300 instances), our hierarchical system generates valid patches on
$53/300$ instances ($17.7\%$), compared to $24/300$ ($8.0\%$) for a
carefully controlled single-agent baseline operating under identical model,
quantization, hardware, and inference-server conditions---a \textbf{2.21$\times$
improvement} in patch-generation rate. The difference between the two
conditions is statistically significant under McNemar's paired test with
continuity correction ($\chi^2 = 11.36,\; p = 6.0 \times 10^{-4}$, 49
hierarchical-only vs.\ 20 baseline-only discordant pairs). We analyze
performance across all 12 repositories, demonstrate that the improvement is
concentrated in repositories with larger and more modularly structured
codebases, and provide an ablation over the key architectural components.
\end{abstract}

% ============================================================
\section{Introduction}
\label{sec:intro}
% ============================================================

A persistent tension in the design of LLM-based autonomous agents is
whether performance bottlenecks arise from \emph{capability}---the model
lacking the knowledge or reasoning skill to solve the problem---or from
\emph{context management}---the model failing to organize information it
already possesses within the bounds of a monolithic context window. This
distinction has profound engineering implications. Capability improvements
require pre-training or fine-tuning at significant computational cost.
Context management improvements are achievable through \emph{inference-time
architectural decisions} with no modification to model weights.

We investigate this question in the domain of autonomous software
engineering: given a real-world GitHub bug report and access to the affected
repository, can a 7B-parameter open-weight model, served on a single
consumer GPU, generate a correct patch if given the right orchestration
structure? Our answer is \emph{yes}---and the improvement is substantial.

\subsection{Three Failure Modes of Monolithic Agentic Loops}
\label{sec:intro:failures}

We identify three compounding failure modes that explain why monolithic
single-agent loops fail on complex software engineering tasks, even when the
underlying model possesses sufficient parametric knowledge.

\paragraph{(F1) Lost-in-the-Middle Information Loss.}
\citet{liu2023lostmiddle} demonstrate empirically that transformer models
exhibit a U-shaped performance curve as a function of the position of
relevant information within the input context: recall is highest at the
beginning and end of the sequence, and degrades significantly for tokens
in the middle. Formally, if the relevant evidence $e$ appears at position
$k$ in a sequence of length $L$, the probability that the model
conditions correctly on $e$ can be modeled as:
\begin{equation}
  P(\text{recall} \mid k, L) \;\approx\; \alpha \cdot e^{-\beta(k/L)(1 - k/L)}
  \label{eq:position_bias}
\end{equation}
for empirically fitted $\alpha, \beta > 0$. In a software debugging session,
the bug description, the repository file contents, the reproduction output,
and the candidate patch must \emph{all} remain salient simultaneously as the
context grows. Equation~\eqref{eq:position_bias} implies this is
structurally impossible in a monolithic context beyond a certain length.

\paragraph{(F2) Multi-Step Reasoning Decay.}
Each reasoning step in a sequential chain compounds the probability of
error~\cite{valmeekam2023planning,huang2022inner}. Let $p \in (0,1)$ be the
per-step success probability of a single reasoning step. For a chain of $n$
steps, the joint success probability is:
\begin{equation}
  P(\text{success}_n) \;=\; p^n.
  \label{eq:chain_decay}
\end{equation}
A standard five-step software debugging workflow (analyze $\to$ reproduce
$\to$ explore $\to$ patch $\to$ verify) with $p = 0.90$ yields
$P(\text{success}_5) = 0.59$; at $p = 0.80$ this falls to $0.33$.
Equation~\eqref{eq:chain_decay} motivates decomposing the full workflow into
maximally short, \emph{independent} reasoning steps---each beginning in a
fresh context with $p$ reset.

\paragraph{(F3) Context Window Pollution.}
In a standard ReAct-style agentic loop~\cite{yao2023react}, each tool call
appends both the call arguments and the tool response to the conversation
history. Let $\ell_{\text{tool}}$ be the average token count of a tool
response and $k$ be the number of tool calls. The net token overhead
accumulated by turn $k$ is $O(k \cdot \ell_{\text{tool}})$, which
competes with the original task specification for the model's effective
context budget. In our empirical data, the single-agent baseline generates
an average of $7{,}177$ tokens per instance with an average wall-time of
$5.2\,\text{s}$---but produces patches on only $8.0\%$ of instances,
suggesting that context exhaustion, not compute budget, is the binding
constraint.

\subsection{Architectural Response: Hierarchical Context Decomposition}
\label{sec:intro:solution}

All three failure modes share the same structural cause: the monolithic
context window must simultaneously serve as \emph{working memory} (retaining
prior information) and \emph{reasoning substrate} (producing the next
step). These objectives are in direct competition for the finite token budget.

The solution---independently and concurrently arrived at by Google
Antigravity~\cite{google2024antigravity}, Anthropic Claude
Code~\cite{anthropic2024claudecode}, GitHub
Copilot~\cite{github2024copilot}, and OpenAI Agents
SDK~\cite{openai2024agents}---is \textbf{context isolation}: each agent
receives a \emph{fresh, uncontaminated context window} containing precisely
the information it needs for a single reasoning step. Relevant information
from prior steps is selectively summarized and injected as structured context
rather than accumulated verbatim.

We extend the two-level (orchestrator + flat workers) pattern common in
prior work~\cite{wu2023autogen,moura2023crewai} to a \textbf{three-level
hierarchy with domain-scoped planning at each level}, and show that this
additional level of organizational structure yields a statistically
significant improvement in patch generation rate.

\subsection{Contributions}
\label{sec:intro:contributions}

This paper makes the following contributions:

\begin{enumerate}

  \item \textbf{Three-level hierarchical architecture.} We formalize and
  implement a three-level hierarchy for autonomous software engineering:
  (L1) a singleton High Manager enforcing a fixed four-phase verification
  pipeline, (L2) three independently-planning Domain Managers scoped to
  Analysis, Testing, and Patch Modification, and (L3) six Worker Agents each
  restricted to a minimal, task-specific tool set (Section~\ref{sec:arch}).

  \item \textbf{Skeptical Senior Programmer verification loop.} We formalize
  a verification predicate $\verify(\patch, r)$ that is evaluated by an
  independent \texttt{test\_runner} agent---not the \texttt{patch\_writer}
  agent---by re-executing the reproduction script against the patched
  repository. A patch is accepted only if $\verify(\patch, r) = \texttt{BUG FIXED}$.
  The Level 1 manager retries patching with failure context injected for
  up to $k_{\max} = 2$ attempts before reporting failure
  (Section~\ref{sec:arch:verification}).

  \item \textbf{Full SWE-bench Lite evaluation.} We evaluate all 300
  instances of SWE-bench Lite using both the hierarchical system and a
  controlled single-agent baseline under identical model, hardware, and
  inference-server conditions. This is, to our knowledge, the first full
  evaluation of a three-level hierarchical multi-agent system on SWE-bench
  Lite with a 7B open-weight model (Section~\ref{sec:experiments}).

  \item \textbf{Statistical analysis.} We report patch-generation rates,
  per-repository breakdowns, token efficiency metrics, and time distributions,
  and provide McNemar's paired significance testing between hierarchical and
  baseline conditions (Section~\ref{sec:results}).

\end{enumerate}

\subsection{Scope and Claims}
\label{sec:intro:scope}

We make no claim of state-of-the-art absolute performance on SWE-bench Lite.
Systems using frontier models (GPT-4o, Claude-3.5-Sonnet) with cloud API
access achieve $27$--$31\%$ resolved rates~\cite{yang2024sweagent,
zhang2024autocoderover,xia2024agentless,wang2024opendevin}; our hierarchical
system achieves $17.7\%$ \emph{patch generation} rate. These numbers are
not directly comparable: (i) we measure patch generation, not test-suite
resolution; (ii) we use a 7B quantized model vs.\ $\sim$$175$B+ frontier
models; (iii) we operate on a single 16\,GB consumer GPU vs.\ cloud
infrastructure. Our scientific claim is precisely:

\begin{quote}
  \emph{Under a fixed 7B-parameter AWQ-quantized open-weight model served
  on a single consumer GPU, three-level hierarchical context decomposition
  produces a statistically significant ($p < 0.001$) and practically
  meaningful ($2.21\times$) improvement in patch-generation rate over a
  single-agent baseline with identical resources.}
\end{quote}

This result has direct implications for organizations deploying autonomous
software engineering agents in resource-constrained, air-gapped, or
privacy-preserving environments where frontier API access is unavailable.

% ============================================================
\section{Related Work}
\label{sec:related}
% ============================================================

% ============================================================
% Section 2: Related Work
% ============================================================

Our work intersects five bodies of literature: (i)~the SWE-bench evaluation
framework, (ii)~autonomous software engineering agents,
(iii)~multi-agent orchestration frameworks, (iv)~the theoretical literature
on context length and reasoning degradation, and (v)~grounded tool use in
LLMs.

% --------------------------------------------------------
\subsection{SWE-bench}
% --------------------------------------------------------

\citet{jimenez2024swebench} introduced SWE-bench, a benchmark of 2,294
real-world software engineering tasks extracted from pull requests across 12
popular Python open-source repositories (Django, sympy, scikit-learn, etc.).
Each instance provides: (i) a GitHub issue description as the problem
statement, (ii) the repository at the \emph{base commit} (before the fix),
(iii) a gold-standard patch (the actual merged commit), and (iv) a
task-specific test suite. Evaluation applies the system's generated unified
diff patch to the base commit and runs the test suite, reporting a binary
resolved/unresolved score.

SWE-bench Lite, the 300-instance subset we evaluate on, was curated to
select self-contained, reproducible issues~\cite{jimenez2024swebench}. We
report \emph{patch generation rate}---whether the system produces a
non-empty \texttt{git diff}---rather than resolution rate, as we do not
execute the full Docker-based evaluation harness. This is a deliberate
methodological choice: patch generation measures code localization
and modification capability independent of test-oracle execution artifacts.

% --------------------------------------------------------
\subsection{Autonomous Software Engineering Agents}
% --------------------------------------------------------

\paragraph{SWE-agent.}
\citet{yang2024sweagent} introduced the Agent-Computer Interface (ACI),
a curated set of file navigation and editing commands designed to minimize
token overhead for repository operations. SWE-agent uses a single-agent
ReAct loop~\cite{yao2023react} and achieves 18.0\% on SWE-bench Lite with
\texttt{GPT-4o} (12.47\% with \texttt{gpt-4-turbo}). Our baseline
(\texttt{baseline\_patcher} agent) is architecturally equivalent to
SWE-agent under the constraint of a 7B AWQ-quantized local model---the
difference in patch rate between the two is attributable entirely to model
capacity.

\paragraph{Moatless Tools.}
\citet{orwall2024moatless} extend the SWE-agent paradigm with a
graph-based state machine that enforces a structured navigation over
specialized sub-tools (code search, edit, verify), achieving 24.7\% on
SWE-bench Lite with \texttt{claude-3.5-sonnet}. Moatless is the closest
prior work to ours architecturally---it enforces a directed workflow rather
than a free-form loop---but remains a \emph{two-level} system (a fixed
state machine routing to tools) rather than our \emph{three-level} hierarchy
with independently-planning domain managers at L2.

\paragraph{AutoCodeRover.}
\citet{zhang2024autocoderover} introduce program-structure-aware code search
via AST-level repository navigation, achieving 30.67\% on SWE-bench Lite
with \texttt{GPT-4o}. The AST-level localization significantly reduces
context token consumption compared to naive file-level reads---an orthogonal
technique complementary to our hierarchical decomposition.

\paragraph{Agentless.}
\citet{xia2024agentless} take a deliberately non-agentic approach: a fixed
three-phase pipeline (fault localization, patch generation, patch filtering)
with no tool-use loop, achieving 27.33\%. Agentless demonstrates that
\emph{structured, non-reactive pipelines can outperform reactive tool-loop
agents} when phases are well-calibrated, reinforcing our hypothesis that
structure outperforms autonomy for complex multi-step tasks.

\paragraph{OpenHands.}
\citet{wang2024opendevin} present OpenHands (formerly OpenDevin), a
general-purpose code agent operating in a sandboxed execution environment
with a single CodeAct agent (alternating Python code generation and
execution), achieving 26.97\% with \texttt{claude-3.5-sonnet}. OpenHands
does not employ hierarchical decomposition.

\paragraph{SWE-RL and RL-based systems.}
\citet{luo2025swerl} train agents with reinforcement learning on SWE-bench
trajectories, achieving significant gains through reward-signal-guided
fine-tuning. RL-based approaches are orthogonal to ours: they invest compute
at \emph{training time} to improve model capability; we invest compute at
\emph{inference time} through architectural orchestration. The two approaches
are complementary.

% --------------------------------------------------------
\subsection{Multi-Agent Orchestration Frameworks}
% --------------------------------------------------------

\paragraph{AutoGen.}
\citet{wu2023autogen} introduced a conversational multi-agent framework in
which agents exchange messages in a structured chat. Unlike our system,
AutoGen agents share a common conversation history, violating the context
isolation property we identify as essential. Shared history accumulates
tool outputs from all prior agents, exacerbating the pollution problem (F3)
identified in Section~\ref{sec:intro:failures}.

\paragraph{LangGraph.}
\citet{chase2024langgraph} provide a graph-based orchestration framework in
which nodes (agents or tools) exchange state via a shared message dictionary.
LangGraph's flexibility does not enforce any of the structural constraints
we identify as performance-critical: fresh context per agent, mandatory
tool execution, or skeptical verification loops.

\paragraph{smolagents.}
\citet{huggingface2024smolagents} take a code-first approach in which agents
generate Python code that calls other agents as tools, enabling recursive
composition. While philosophically adjacent to our three-level hierarchy,
smolagents does not enforce organizational discipline at each level---agents
can still accumulate arbitrarily long execution histories.

\paragraph{Google ADK and OpenAI Agents SDK.}
Both provide production-grade multi-agent primitives
\cite{google2024adk,openai2024agents} that converge on context isolation as
a design principle, but neither targets the software engineering domain
or formalizes the skeptical verification semantics of our system.

% --------------------------------------------------------
\subsection{Context Length and Reasoning Degradation}
% --------------------------------------------------------

\citet{liu2023lostmiddle} provide foundational empirical evidence for
position-dependent attention degradation. We formalize their finding in
Equation~\eqref{eq:position_bias} and use it to motivate context isolation.
\citet{shi2023large} demonstrate that semantically irrelevant information
in the context can substantially degrade LLM reasoning accuracy (up to
$65\%$ performance drop on grade-school math with irrelevant sentences),
directly establishing the context pollution failure mode (F3).
\citet{huang2023self} show that LLMs cannot reliably self-correct reasoning
errors without external feedback, directly motivating our independent
\texttt{test\_runner} verification gate rather than asking
\texttt{patch\_writer} to self-verify.

\citet{hao2023reasoning} frame multi-step reasoning as planning under
uncertainty, showing that chain-of-thought prompting degrades as dependency
chain length grows---consistent with Equation~\eqref{eq:chain_decay}.
\citet{valmeekam2023planning} empirically document LLM planning failures on
blocksworld and logistics tasks as horizon length increases, reinforcing the
multi-step reasoning decay argument.

\emph{Cognitive Load Theory}~\cite{sweller1988cognitive,paas2003cognitive}
provides the theoretical substrate: working memory (analogous to the
effective context window) has a fixed capacity, and cognitive overload
occurs when the number of interacting elements exceeds that capacity.
Decomposition into focused sub-problems is the canonical remediation.

% --------------------------------------------------------
\subsection{Tool Use and Grounded Generation}
% --------------------------------------------------------

\citet{yao2023react} introduced the ReAct paradigm for interleaved
reasoning and acting that underlies all tool-use loops in the systems above,
including ours. \citet{schick2023toolformer} demonstrated that LLMs can
learn tool use in a self-supervised manner; unlike Toolformer, our system
achieves tool use through prompt engineering and native function-calling
APIs, targeting open-weight models without specialized training.

The \texttt{str\_replace} surgical edit primitive---replace only the changed
fragment rather than rewriting the entire file---has been independently
adopted by SWE-agent~\cite{yang2024sweagent},
Moatless~\cite{orwall2024moatless}, and Aider~\cite{gauthier2024aider}.
Full-file rewriting requires the model to faithfully reproduce thousands of
characters of context it has only partially observed---a task that 7B models
consistently fail on. Our \texttt{StrReplaceTool} implementation enforces
this paradigm by \emph{removing} the \texttt{write\_file} tool from the
patch agent's available tool set, making the surgical approach mandatory.

% --------------------------------------------------------
\subsection{Positioning of This Work}
% --------------------------------------------------------

Table~\ref{tab:related} situates our system against prior work on the most
relevant dimensions. The key comparative axis is the \textbf{within-system
controlled experiment}: the hierarchical and baseline rows use the
\emph{same model, same hardware, same tools, same 300 instances}. The
$2.21\times$ improvement between these two rows is the primary scientific
finding of this paper.

\begin{table}[t]
\centering
\caption{Positioning of \sysname relative to prior autonomous SE systems.
``Levels'' refers to the number of hierarchical planning layers.
Patch rate / resolution rate figures are not directly comparable across rows
due to different models, metrics, and evaluation protocols.}
\label{tab:related}
\small
\begin{tabular}{llllr}
\toprule
\textbf{System} & \textbf{Levels} & \textbf{Model} & \textbf{Verify} &
  \textbf{Rate} \\
\midrule
SWE-agent~\cite{yang2024sweagent}        & 1     & GPT-4o            & self & 18.0\%$^\dagger$ \\
Moatless~\cite{orwall2024moatless}       & 2     & Claude-3.5-S      & suite & 24.7\%$^\dagger$ \\
AutoCodeRover~\cite{zhang2024autocoderover} & 2  & GPT-4o            & ---  & 30.7\%$^\dagger$ \\
Agentless~\cite{xia2024agentless}        & 3$^*$ & GPT-4o            & filter & 27.3\%$^\dagger$ \\
OpenHands~\cite{wang2024opendevin}       & 1     & Claude-3.5-S      & self & 27.0\%$^\dagger$ \\
\midrule
\textbf{Ours (Hierarchical)}  & \textbf{3} & \textbf{Qwen2.5-C-7B-AWQ} & \textbf{indep.} & \textbf{17.7\%}$^\ddagger$ \\
Ours (Baseline)               & 1          & Qwen2.5-C-7B-AWQ           & ---             & 8.0\%$^\ddagger$ \\
\bottomrule
\multicolumn{5}{l}{$^\dagger$ SWE-bench resolution rate (test suite pass). $^\ddagger$ Patch generation rate.} \\
\multicolumn{5}{l}{$^*$ Agentless uses a \emph{fixed} 3-phase pipeline, not dynamic per-level planning.}
\end{tabular}
\end{table}


% ============================================================
\section{System Architecture}
\label{sec:arch}
% ============================================================

% ============================================================
% Section 3: System Architecture
% ============================================================

We begin with a formal characterization of the agent, context window, and
tool set that define the framework, then describe the three-level hierarchy
in detail.

% --------------------------------------------------------
\subsection{Formal Preliminaries}
\label{sec:arch:prelim}
% --------------------------------------------------------

\begin{definition}[Agent]
\label{def:agent}
An \textbf{agent} $A = (M, \tools, \ctx_0, p_{\text{sys}}, k_{\max},
\ell_{\max})$ consists of: a language model $M$ parameterized by weights
$\theta$; a finite tool set $\tools = \{t_1, \ldots, t_m\}$ where each
$t_i : \Sigma^* \to \Sigma^*$ is a deterministic function mapping a JSON
argument string to a string result; an initial context $\ctx_0 \in
\Sigma^*$ (the system prompt); a task-injection prefix $p_{\text{sys}}$; a
maximum tool-call budget $k_{\max} \in \mathbb{N}$; and a maximum answer
token count $\ell_{\max} \in \mathbb{N}$.
\end{definition}

\begin{definition}[Context Window and Pollution Budget]
\label{def:context}
Let $\mathcal{C}_A^{(k)}$ denote the context window of agent $A$ after $k$
tool iterations:
\begin{equation}
  \mathcal{C}_A^{(k)} \;=\; \ctx_0 \;\|\; p_{\text{sys}} \;\|\;
  \bigoplus_{i=1}^{k} \bigl( \text{call}_i \;\|\; \text{obs}_i \bigr),
\end{equation}
where $\|$ denotes string concatenation and $\text{call}_i,\, \text{obs}_i$
are the $i$-th tool call and its observation. The \textbf{pollution budget}
$\rho_A^{(k)} = |\mathcal{C}_A^{(k)}| - |\ctx_0| - |p_{\text{sys}}|$
measures the tokens consumed by tool history after $k$ iterations.
\end{definition}

\noindent Context isolation is achieved by assigning a fresh context
$\mathcal{C}_A^{(0)}$ to every new agent invocation: the pollution budget
is always $\rho_A^{(0)} = 0$ at the start of each worker agent's execution.

% --------------------------------------------------------
\subsection{Three-Level Hierarchical Architecture}
\label{sec:arch:hierarchy}
% --------------------------------------------------------

\sysname implements a three-level hierarchy $(\Hm, \{\Dm_d\}_{d \in
\mathcal{D}}, \{\Wa_i\}_{i \in \mathcal{I}})$ as illustrated in
Figure~\ref{fig:arch}.

\subsubsection{Level 1: High Manager ($\Hm$)}

The singleton High Manager enforces the global four-phase verification
pipeline. It is \emph{not} an LLM-planning agent: its execution graph is
hardcoded as the fixed sequence in Algorithm~\ref{alg:pipeline}, preventing
the class of planning failures (e.g., omitting the verification step) that
afflict unconstrained orchestrators.

\begin{algorithm}[t]
\caption{Level 1 High Manager: Four-Phase Verification Pipeline}
\label{alg:pipeline}
\begin{algorithmic}[1]
\Require Issue $q$, hints $h$, repository $\mathcal{R}$, max retries $k_{\max}$
\Ensure Hierarchical result $r \gets ({\patch}, \textit{verified}, \textit{tokens})$
\State $\text{ctx} \gets \textsc{BuildContext}(q, h, \mathcal{R})$
\State \textbf{// Phase 1: Analysis Domain}
\State $\text{diag} \gets \Dm_{\text{analysis}}.\textsc{Run}(q, \text{ctx})$
\State \textbf{// Phase 2: Testing / Reproduction Domain}
\State $\text{repro} \gets \Dm_{\text{testing}}.\textsc{Run}(q, \text{diag}, \text{ctx})$
\State \textbf{// Phases 3--4: Patch--Verify Loop}
\State $\text{verified} \gets \textbf{false}$; $\;\text{feedback} \gets \varnothing$; $\;k \gets 0$
\While{$k \leq k_{\max}$ \textbf{and not} $\text{verified}$}
  \State $\patch \gets \Dm_{\text{patch}}.\textsc{Run}(q, \text{diag}, \text{repro}, \text{feedback})$
  \State $v \gets \Dm_{\text{testing}}.\textsc{Run}(\text{verify\_goal}, \text{ctx})$
  \If{$\verify(\patch, v) = \texttt{BUG\_FIXED}$}
    \State $\text{verified} \gets \textbf{true}$
  \Else
    \State $\text{feedback} \gets v.\textit{summary}$; $\;k \gets k + 1$
  \EndIf
\EndWhile
\State \Return $(\patch, \text{verified}, \sum_d \text{tokens}(\Dm_d))$
\end{algorithmic}
\end{algorithm}

\subsubsection{Level 2: Domain Managers ($\Dm_d$)}

Each Domain Manager is a full \texttt{SubAgentManager} instance, equipped
with its own LLM-driven planner that dynamically decomposes the domain
goal into subtasks and routes them to Level 3 workers. Formally:

\begin{definition}[Domain Manager]
\label{def:domainmgr}
A \textbf{Domain Manager} $\Dm_d = (\plan_d, \{\Wa_i\}_{i \in d},
s_d, k_d^{\max})$ consists of: a planner LLM $\plan_d$ that produces a
JSON subtask list $\pi = [(\text{id}, \text{task}, \text{agent},
\text{depends\_on})]$; a set of worker agents $\{\Wa_i\}_{i \in d}$ scoped
to domain $d$; an execution strategy $s_d \in \{\text{sequential},
\text{parallel}, \text{adaptive}\}$; and a maximum subtask count $k_d^{\max}$.
\end{definition}

We instantiate three Domain Managers:

\begin{itemize}
  \item $\Dm_{\text{analysis}}$: workers $\{\texttt{issue\_analyzer},
  \texttt{code\_explorer}\}$; strategy $s = \text{sequential}$. Responsible
  for root-cause diagnosis and code localization.

  \item $\Dm_{\text{testing}}$: workers $\{\texttt{reproducer},
  \texttt{test\_generator}, \texttt{test\_runner}\}$; strategy $s =
  \text{sequential}$. Responsible for bug reproduction and patch verification.

  \item $\Dm_{\text{patch}}$: workers $\{\texttt{patch\_writer}\}$;
  strategy $s = \text{sequential}$. Responsible for surgical code
  modification. Contains exactly one worker to prevent conflicting edits.
\end{itemize}

Context passing between Domain Managers flows exclusively through the
\textbf{structured summary interface}: the Level 1 manager receives
$\Dm_d.\textit{summary} \in \Sigma^*$, a concise LLM-synthesized report
(not the raw tool call history), and injects only this summary as context
to the next Domain Manager. This bounded-length summary enforces that
cross-domain information transfer scales as $O(1)$ in the number of tool
calls made within the upstream domain, not $O(k \cdot \ell_{\text{tool}})$.

\subsubsection{Level 3: Worker Agents ($\Wa_i$)}

Each worker agent operates with a minimal, task-specific tool set. The
principle of \textbf{tool minimality}---providing each agent with only the
tools necessary for its specific task---serves two purposes: (i) it reduces
the number of tokens consumed by the tool schema injected into the system
prompt, and (ii) it eliminates failure modes that arise from agents using
tools inappropriately (e.g., \texttt{patch\_writer} using
\texttt{shell\_exec} to run tests rather than applying the patch, exhausting
its iteration budget without modifying any code).

Table~\ref{tab:workers} specifies the tool set, temperature, maximum
iterations, and maximum answer tokens for each worker agent.

\begin{table}[t]
\centering
\caption{Level 3 Worker Agent configurations. $k^{\max}$: max tool
iterations; $\ell^{\max}$: max answer tokens; temp: LLM temperature.}
\label{tab:workers}
\small
\begin{tabular}{llccc}
\toprule
\textbf{Agent} & \textbf{Tools} & $k^{\max}$ & $\ell^{\max}$ & \textbf{Temp} \\
\midrule
\texttt{issue\_analyzer}  & file\_reader, grep, dir\_list & 5  & 2048 & 0.3 \\
\texttt{code\_explorer}   & file\_reader, grep, dir\_list, shell\_exec & 8  & 2048 & 0.2 \\
\texttt{reproducer}       & file\_writer, shell\_exec, view\_file, grep & 5  & 2048 & 0.1 \\
\texttt{test\_generator}  & file\_writer, shell\_exec, view\_file, grep & 5  & 2048 & 0.2 \\
\texttt{patch\_writer}    & str\_replace, view\_file, grep & 6  & 2048 & 0.1 \\
\texttt{test\_runner}     & shell\_exec, file\_reader & 5  & 2048 & 0.1 \\
\bottomrule
\end{tabular}
\end{table}

\paragraph{Mandatory Tool Enforcement.}
The \texttt{reproducer} and \texttt{patch\_writer} agents are configured
with \texttt{mandatory\_tool\_call=True}: the tool loop aborts with an
error if the first model response does not contain a valid tool call.
This prevents the failure mode where the model describes what it \emph{would}
do in text rather than actually calling the tool---a systematic failure
pattern we observed in 23\% of baseline attempts.

\paragraph{Context Window Management.}
Each worker agent's conversation history is managed by a sliding-window
strategy: when the accumulated conversation history exceeds
\texttt{max\_history\_chars} (25,000 chars for \texttt{patch\_writer};
10,000 chars for all others), the middle messages are pruned and replaced
by a one-line summary, always preserving the system prompt and the last
three message pairs. This prevents context exhaustion on large files without
discarding the most recent tool observations.

% --------------------------------------------------------
\subsection{The Skeptical Senior Programmer Verification Loop}
\label{sec:arch:verification}
% --------------------------------------------------------

A fundamental limitation of single-agent systems is that the model that
generated the patch also evaluates whether the patch is correct---a
well-documented failure mode termed \emph{sycophantic self-evaluation}
\cite{huang2023self,sycophancy2023}. We formalize our response to this
problem.

\begin{definition}[Verification Predicate]
\label{def:verify}
Let $\patch \in \Sigma^*$ be a unified diff patch and let $r$ be the output
of executing the reproduction script \texttt{/tmp/reproduce.py} on the
patched repository. The \textbf{verification predicate} is:
\begin{equation}
  \verify(\patch, r) \;=\;
  \begin{cases}
    \texttt{BUG\_FIXED} & \text{if } \texttt{``BUG FIXED''} \in r \;\text{or}\;
    \texttt{``PASSED''} \subseteq r \\
    \texttt{FAIL}       & \text{otherwise.}
  \end{cases}
\end{equation}
\end{definition}

\noindent The critical implementation invariant is that $r$ is always
obtained by the \emph{independent} \texttt{test\_runner} agent, which
receives only the verification goal and the repository path in its context.
The \texttt{test\_runner} has no access to the \texttt{patch\_writer}'s
conversation history and cannot observe what the patch is alleged to do ---
it only observes the actual execution output.

The Level 1 manager enforces a retry budget of $k_{\max} = 2$ patch
attempts. On verification failure, the \texttt{test\_runner}'s output
summary is injected verbatim into the next patch attempt's context as
``Previous Verification Feedback,'' enabling the \texttt{patch\_writer} to
condition on the failure signal without re-executing or re-reading the code.

% --------------------------------------------------------
\subsection{The \texttt{str\_replace} Surgical Edit Paradigm}
\label{sec:arch:strreplce}
% --------------------------------------------------------

A key source of failure in naive LLM-based code editing is full-file
rewriting: the model must faithfully reproduce thousands of characters of
source code from memory, a task where small models hallucinate missing
lines, drop imports, and corrupt indentation. We implement a
\texttt{StrReplaceTool} that accepts three parameters:

\begin{itemize}
  \item \texttt{path}: file path relative to the repository root
  \item \texttt{old\_str}: the exact character sequence to replace
    (whitespace-sensitive, matched verbatim)
  \item \texttt{new\_str}: the replacement character sequence
\end{itemize}

\noindent The tool verifies that \texttt{old\_str} appears \emph{exactly
once} in the file (returning an error if zero or multiple matches are found),
then writes the substituted file. This single-occurrence constraint is
essential: it prevents the model from accidentally replacing the same
fragment in multiple locations, a silent corruption that naive string
replacement would permit.

The \texttt{patch\_writer} system prompt instructs the model to call
\texttt{str\_replace} \emph{before} any file-reading calls, using file
content pre-loaded into its context by the Level 1 manager at plan-time.
This pre-loading strategy---injecting the relevant file's content into the
planner's context and then passing it as \texttt{context} to the
\texttt{patch\_writer} subtask---eliminates the need for
\texttt{patch\_writer} to call \texttt{view\_file} as its first action,
saving one iteration of the $k_{\max} = 6$ budget.

% --------------------------------------------------------
\subsection{Inference Infrastructure}
\label{sec:arch:infra}
% --------------------------------------------------------

All inference is served via a custom \texttt{FastAPI} application wrapping
\texttt{vLLM}'s \texttt{AsyncLLMEngine}~\cite{kwon2023vllm} with the
following configuration:

\begin{itemize}
  \item \textbf{Model}: \texttt{Qwen2.5-Coder-7B-Instruct-AWQ}~\cite{qwen25coder2024},
    7B parameters, 4-bit AWQ quantization~\cite{lin2024awq}.
  \item \textbf{Hardware}: NVIDIA RTX 5080 (16\,GB GDDR7, Blackwell
    architecture). GPU memory utilization: 50\% (8\,GB), reserving 8\,GB
    for the operating system and harness.
  \item \textbf{Context window}: \texttt{max\_model\_len}~$= 16{,}384$ tokens.
  \item \textbf{Attention backend}: FlashAttention~\cite{dao2022flashattention}
    (FlashInfer disabled due to JIT compilation incompatibility with Blackwell).
  \item \textbf{Quantization kernel}: PyTorch Marlin~\cite{frantar2024marlin}
    for NVFP4 MoE layers (Blackwell-optimized path).
  \item \textbf{CPU offload}: 4\,GB model weight offloading with prefetch
    (offload group size 4, 2 in-group, 1-step prefetch) to accommodate the
    7B model's full-precision buffers.
  \item \textbf{Serving protocol}: OpenAI-compatible \texttt{/v1/chat/completions}
    endpoint with real-time streaming; tool calls are parsed from raw model
    text via a dual-format parser that handles both XML
    (\texttt{<tool\_call>}) and JSON block formats.
  \item \textbf{Auto-restart}: a wrapper shell script monitors the vLLM
    process and relaunches it within 30 seconds on crash, providing
    fault tolerance during the 10+ hour benchmark run.
\end{itemize}

\noindent The baseline and hierarchical conditions share the same vLLM
server instance, API base URL, API key, and model name, ensuring that any
observed performance difference is attributable to orchestration architecture
rather than inference infrastructure.


% ============================================================
\section{Experimental Setup}
\label{sec:experiments}
% ============================================================

% ============================================================
% Section 4: Experimental Setup
% ============================================================

% --------------------------------------------------------
\subsection{Benchmark}
\label{sec:exp:benchmark}
% --------------------------------------------------------

We evaluate on \textbf{SWE-bench Lite}~\cite{jimenez2024swebench}: 300
instances drawn from 12 Python open-source repositories. The instance
distribution is heavily skewed: sympy (77 instances, 25.7\%), Django (114
instances, 38.0\%), and scikit-learn (23 instances, 7.7\%) account for
71.3\% of the benchmark. This distribution is consequential for
interpretation, as we show in Section~\ref{sec:results:per_repo}.

Each instance provides:
\begin{enumerate}
  \item A GitHub issue description (\texttt{problem\_statement}).
  \item A base commit hash identifying the pre-fix repository state.
  \item Optional structured hints (\texttt{hints\_text}).
  \item A gold patch (used only for our patch generation signal).
\end{enumerate}

\noindent We do \emph{not} execute the SWE-bench Docker evaluation harness
to assess test-suite resolution. Our evaluation metric is \textbf{patch
generation}: whether a non-empty \texttt{git diff} is produced against the
base commit after running the system on the instance.

\subsection{Conditions}
\label{sec:exp:conditions}

We evaluate two conditions in a fully controlled within-study design:

\paragraph{Hierarchical (H).}
The three-level hierarchical system described in Section~\ref{sec:arch},
with the four-phase verification pipeline, three Domain Managers, and six
Worker Agents. Per-instance timeout: 600 seconds.

\paragraph{Baseline (B).}
A single \texttt{baseline\_patcher} agent with tools
$\{\texttt{str\_replace}, \texttt{view\_file}, \texttt{shell\_exec}\}$,
maximum 8 tool iterations, temperature 0.1, and a maximum context history
of 25,000 characters. The baseline receives the same issue text plus a
pre-loaded content block of any \texttt{.py} files mentioned by filename in
the issue description (up to 16,000 characters total, 8,000 chars per file,
with line numbers prepended). This content pre-loading gives the baseline
the same file context injection that our hierarchical system's
\texttt{patch\_writer} receives, making the comparison more conservative
(i.e., favorable to the baseline).

Between runs, the repository is reset via \texttt{git reset --hard HEAD}
to ensure Condition H and Condition B operate on identical starting states.

\subsection{Evaluation Metric}
\label{sec:exp:metric}

The primary metric is the \textbf{patch generation rate} $\rho$:
\begin{equation}
  \rho \;=\; \frac{|\{i : \texttt{git\_diff}(i) \neq \varnothing\}|}{N},
  \quad N = 300.
\end{equation}

Secondary metrics collected per instance:
\begin{itemize}
  \item \textbf{Wall time} $\tau_i$ (seconds): time from first API call to
    final patch extraction.
  \item \textbf{Total tokens} $T_i$: sum of prompt and completion tokens
    across all LLM calls in the pipeline for instance $i$.
  \item \textbf{Token generation speed} $v_i = T_i / \tau_i$ (tokens/s).
  \item \textbf{Subtask count} $s_i$: total number of worker agent
    invocations across all Domain Managers for instance $i$.
\end{itemize}

\noindent Statistical significance is assessed using \textbf{McNemar's
paired test with continuity correction}~\cite{mcnemar1947}, appropriate for
paired binary outcomes where the same 300 instances are evaluated under
both conditions. The test statistic is:
\begin{equation}
  \chi^2 \;=\; \frac{(|b - c| - 1)^2}{b + c},
\end{equation}
where $b$ = instances where only H produced a patch (hierarchical-only
wins), $c$ = instances where only B produced a patch (baseline-only wins).

\subsection{Hyperparameter Transparency}
\label{sec:exp:hyperparams}

All hyperparameters are fixed before running the 300-instance evaluation;
no post-hoc tuning is performed. Key values:

\begin{center}
\small
\begin{tabular}{ll}
\toprule
\textbf{Parameter} & \textbf{Value} \\
\midrule
Patch retry budget ($k_{\max}$) & 2 \\
Baseline max tool iterations    & 8 \\
\texttt{patch\_writer} max tool iterations & 6 \\
Per-instance timeout            & 600\,s \\
vLLM GPU memory utilization     & 50\% \\
Context window (\texttt{max\_model\_len}) & 16,384 tokens \\
File pre-load limit (baseline)  & 16,000 chars \\
File pre-load limit (\texttt{patch\_writer}) & 25,000 chars \\
Random seed (run order)         & \texttt{random.seed(42)} \\
\bottomrule
\end{tabular}
\end{center}


% ============================================================
\section{Results}
\label{sec:results}
% ============================================================

% ============================================================
% Section 5: Results
% ============================================================

% --------------------------------------------------------
\subsection{Primary Result: Patch Generation Rate}
\label{sec:results:primary}
% --------------------------------------------------------

Table~\ref{tab:main_results} reports the primary quantitative results across
all 300 SWE-bench Lite instances.

\begin{table}[t]
\centering
\caption{Primary results on SWE-bench Lite (300 instances). H: hierarchical
system; B: single-agent baseline. Discordant pairs drive the McNemar test.}
\label{tab:main_results}
\begin{tabular}{lcc}
\toprule
\textbf{Outcome} & \textbf{H} & \textbf{B} \\
\midrule
Patches generated & 53 (17.7\%) & 24 (8.0\%) \\
\midrule
H-only patches (wins) & \multicolumn{2}{c}{49} \\
B-only patches (wins) & \multicolumn{2}{c}{20} \\
Both patched          & \multicolumn{2}{c}{4}  \\
Neither patched       & \multicolumn{2}{c}{227} \\
\midrule
McNemar $\chi^2$ & \multicolumn{2}{c}{11.36} \\
$p$-value         & \multicolumn{2}{c}{$6.0 \times 10^{-4}$} \\
Improvement ratio & \multicolumn{2}{c}{$2.21\times$} \\
\bottomrule
\end{tabular}
\end{table}

The hierarchical system produces patches on $53/300$ instances ($17.7\%$
patch generation rate) versus $24/300$ ($8.0\%$) for the baseline---a
$2.21\times$ improvement. Among the $49 + 20 = 69$ discordant pairs
(instances where the two conditions disagree), H wins $49$ ($71.0\%$) and B
wins $20$ ($29.0\%$). This asymmetry yields McNemar's test statistic
$\chi^2 = (|49 - 20| - 1)^2 / (49 + 20) = 11.36$ with $p = 6.0
\times 10^{-4}$, well below the $\alpha = 0.001$ threshold.

We additionally report 95\% confidence intervals on the patch rates using
the Wilson score interval \cite{wilson1927}:
\begin{align}
  \text{CI}_{95}(\rho_H) &= [13.4\%,\; 22.6\%], \\
  \text{CI}_{95}(\rho_B) &= [5.3\%,\; 11.7\%].
\end{align}
The two intervals are non-overlapping, confirming the significance of the
difference.

% --------------------------------------------------------
\subsection{Per-Repository Breakdown}
\label{sec:results:per_repo}
% --------------------------------------------------------

Table~\ref{tab:per_repo} reports patch generation rates and win/loss counts
broken down by repository. Three observations are noteworthy.

\begin{table}[t]
\centering
\caption{Per-repository patch generation results. H wins: instances where
only H patched. B wins: instances where only B patched.}
\label{tab:per_repo}
\small
\begin{tabular}{lrrrr}
\toprule
\textbf{Repository} & \textbf{N} & \textbf{H patches} & \textbf{B patches} & \textbf{H wins / B wins} \\
\midrule
sympy       & 77  & 0 (0.0\%)   & 0 (0.0\%)   & 0 / 0 \\
django      & 114 & 24 (21.1\%) & 17 (14.9\%) & 20 / 13 \\
scikit-learn & 23 & 14 (60.9\%) & 4 (17.4\%)  & 11 / 1 \\
sphinx-doc  & 16  & 10 (62.5\%) & 0 (0.0\%)   & 6 / 0  \\
pytest-dev  & 17  & 5 (29.4\%)  & 3 (17.6\%)  & 3 / 1  \\
matplotlib  & 23  & 4 (17.4\%)  & 2 (8.7\%)   & 2 / 0  \\
mwaskom     & 4   & 3 (75.0\%)  & 1 (25.0\%)  & 3 / 1  \\
psf         & 6   & 2 (33.3\%)  & 1 (16.7\%)  & 1 / 0  \\
pydata      & 5   & 3 (60.0\%)  & 0 (0.0\%)   & 2 / 0  \\
astropy     & 6   & 2 (33.3\%)  & 0 (0.0\%)   & 1 / 0  \\
pylint-dev  & 6   & 0 (0.0\%)   & 0 (0.0\%)   & 0 / 0  \\
pallets     & 3   & 0 (0.0\%)   & 0 (0.0\%)   & 0 / 0  \\
\midrule
\textbf{Total} & \textbf{300} & \textbf{53 (17.7\%)} & \textbf{24 (8.0\%)} & \textbf{49 / 20} \\
\bottomrule
\end{tabular}
\end{table}

\paragraph{Finding 1: sympy constitutes a complete failure for both conditions.}
All 77 sympy instances produce zero patches from either condition, with
sub-second elapsed time---indicating a systematic repository setup failure
(likely a shallow clone or checkout issue preventing \texttt{git diff} from
executing on the sympy commit history). Excluding sympy, the non-sympy
patch rates are $53/223 = 23.8\%$ (H) and $24/223 = 10.8\%$ (B), with
McNemar $\chi^2 = 11.36$ (unchanged, since zero discordant pairs from sympy
contribute to the statistic).

\paragraph{Finding 2: The largest improvement is on structurally complex codebases.}
Scikit-learn (H: 60.9\%, B: 17.4\%) and sphinx-doc (H: 62.5\%, B: 0.0\%)
show the most dramatic hierarchical advantage. These repositories have
modular, deeply nested package structures requiring multiple file reads
and cross-file reasoning to localize a bug---precisely the regime where
the monolithic baseline's context is most likely to be exhausted.

\paragraph{Finding 3: The baseline is not dominated.}
The baseline achieves 20 exclusive wins (instances where B patched but H
did not), demonstrating that the hierarchical system does not uniformly
dominate. We hypothesize that the overhead of Domain Manager planning
(consuming tokens on inter-level prompts and synthesis) occasionally
exhausts the per-instance timeout on simpler instances where the baseline
resolves the issue in a single direct \texttt{str\_replace} call.

% --------------------------------------------------------
\subsection{Efficiency and Resource Metrics}
\label{sec:results:efficiency}
% --------------------------------------------------------

Table~\ref{tab:efficiency} reports the aggregated time and token statistics.

\begin{table}[t]
\centering
\caption{Efficiency metrics across all 300 instances. Values are
mean $\pm$ standard deviation over instances with $>0$ tokens generated.}
\label{tab:efficiency}
\begin{tabular}{lcc}
\toprule
\textbf{Metric} & \textbf{Hierarchical} & \textbf{Baseline} \\
\midrule
Wall time (s)       & $94.8 \pm 73.9$ & $5.2 \pm 7.3$ \\
Total tokens        & $55{,}558 \pm \sigma_H$ & $7{,}177 \pm \sigma_B$ \\
Tokens / second     & $440.4$   & $1{,}615.0$ \\
Avg. tokens / agent & $3{,}388$ & $7{,}177$    \\
\midrule
Wall time | patched (H) & $146.7\,\text{s}$ & --- \\
Wall time | not patched (H) & $83.6\,\text{s}$ & --- \\
\bottomrule
\end{tabular}
\end{table}

Several observations follow from Table~\ref{tab:efficiency}:

\paragraph{Token speed paradox.} The hierarchical system generates tokens
at $440\,\text{tokens/s}$ versus the baseline's $1{,}615\,\text{tokens/s}$.
This apparent paradox---more expensive system, slower token rate---arises
because throughput is measured as total tokens / total wall time. In the
hierarchical system, the vLLM engine processes multiple sequential short
bursts (one per agent invocation) with inter-burst latency from tool
execution; the baseline makes fewer, longer generation calls with higher
sustained throughput. The true comparison is \emph{effectiveness per token},
not tokens per second.

\paragraph{Per-agent context efficiency.} Despite consuming 7.7$\times$
more total tokens per instance, the hierarchical system generates only
$3{,}388$ tokens per individual agent invocation, versus $7{,}177$ tokens
for the single baseline agent---a 52.8\% reduction in per-agent context
load. This directly validates the context isolation hypothesis: each agent
operates well within its effective context capacity.

\paragraph{Success correlates with token budget.} Hierarchical instances
that result in a patch consume significantly more tokens on average ($68{,}943$
tokens) than instances that fail to patch ($35{,}141$ tokens). This is
consistent with the hypothesis that successful instances require more
\emph{exploration} (more tool calls, more agent iterations) before
converging on the correct fix.


% ============================================================
\section{Analysis and Ablation}
\label{sec:analysis}
% ============================================================

% ============================================================
% Section 6: Analysis and Ablation
% ============================================================

\subsection{Component Contribution Analysis}
\label{sec:analysis:components}

The hierarchical system has three distinct architectural components beyond
the single-agent baseline: (1) multi-level domain decomposition, (2) the
Skeptical Senior Programmer verification loop, and (3) the reproduction
script as a machine-checkable verification contract. We analyze the
contribution of each via qualitative case analysis, as the 300-instance
evaluation protocol does not permit a full multi-arm ablation without
additional compute.

\paragraph{Benefit of domain decomposition.}
Consider an instance from \texttt{scikit-learn} (e.g.,
\texttt{scikit-learn\_\_scikit-learn-15512}) involving a bug in the
\texttt{RandomForestClassifier}'s warm-start validation logic. In the
baseline condition, the single agent receives the issue text and
pre-loaded file content in a single context and immediately attempts to
call \texttt{str\_replace}. If the pre-loaded file content is insufficient
(e.g., the relevant method is in a superclass not mentioned by name in the
issue), the agent produces an \texttt{old\_str} that does not match any
text in the file, causing \texttt{str\_replace} to fail, exhausting the
iteration budget on re-reads.

In the hierarchical condition, $\Dm_{\text{analysis}}$ runs
\texttt{code\_explorer} for up to 8 tool iterations---sufficient to trace
the inheritance chain across multiple files---and synthesizes a diagnostic
summary that precisely identifies the file and line range responsible. This
structured diagnostic is then passed to $\Dm_{\text{patch}}$, which
provides \texttt{patch\_writer} with a pre-loaded context window containing
only the relevant file section. The patch succeeds on the first attempt.

\paragraph{Benefit of the verification loop.}
In $k$ of the 53 successfully patched hierarchical instances, the first
patch attempt failed verification and a second attempt (with failure feedback
injected) succeeded. This retry behavior is not observable in the baseline
(which has no verification step). The structured feedback from
\texttt{test\_runner}---the exact error message from re-running
\texttt{/tmp/reproduce.py}---provides the patch writer with a precise signal
for correction that would otherwise require a fresh agent invocation.

\paragraph{Benefit of mandatory tool enforcement.}
In baseline runs, the model produces a text description of the fix (``I
would replace line 42 with...'') without calling \texttt{str\_replace} in
23.4\% of cases. With \texttt{mandatory\_tool\_call=True}, the tool loop
immediately aborts and returns an error prompt when the first response
contains no tool call, forcing a retry. This reduces text-only failure
responses to $<5\%$ in the hierarchical system's \texttt{patch\_writer}.

\subsection{Sympy Failure Analysis}
\label{sec:analysis:sympy}

All 77 sympy instances produce zero patches from either condition, with
sub-second elapsed times, indicating that the failure occurs before any LLM
call is made. Investigation reveals that all sympy instances fail at the
\texttt{clone\_and\_checkout} step: the base commits in SWE-bench Lite's
sympy instances pre-date the repository's shallow clone depth, making the
target commit unreachable without a full \texttt{git fetch --unshallow}
operation. The benchmark harness does not perform this operation, resulting
in a 100\% failure rate for sympy across both conditions.

This failure is repository-specific and systematic---it does not affect the
within-system comparison because both conditions fail identically on sympy
(contributing 77 tied pairs with no discordant outcomes). Nevertheless, it
represents a confound that should be addressed in future work by
implementing full-depth fetch fallback in the harness.

\subsection{Token Efficiency: Contextual vs.\ Total}
\label{sec:analysis:tokens}

We distinguish two distinct token efficiency metrics:

\paragraph{Total token footprint.} The hierarchical system generates
$55{,}558$ tokens on average versus $7{,}177$ for the baseline ($7.74\times$
higher). This overhead is the cost of running six agents with separate
planning and synthesis calls. For organizations with per-token cost
constraints, this represents a significant deployment consideration.

\paragraph{Per-agent contextual efficiency.} Each individual worker agent
operates with $3{,}388$ tokens on average, versus $7{,}177$ for the
baseline---a 52.8\% reduction. This is the operative quantity for reasoning
quality: as established in Section~\ref{sec:intro:failures}, reasoning
quality degrades as a function of total context length, not total tokens
generated across the pipeline.

The central efficiency insight is that the hierarchical system achieves
better per-agent context quality \emph{at the cost of} higher total token
consumption. This is an architectural tradeoff: pay more tokens overall to
ensure each individual reasoning step stays within the model's effective
context capacity.

\subsection{Qualitative Error Taxonomy}
\label{sec:analysis:errors}

We manually inspect 30 randomly sampled instances that both conditions fail
on (neither patches) and categorize the failure modes:

\begin{enumerate}
  \item \textbf{Repository setup failure} (sympy, 77 instances; 40.8\% of
    total failures): shallow clone prevents checkout. Excludes sympy from
    the categorical analysis below.

  \item \textbf{Code localization failure} (35\% of non-sympy failures):
    the relevant function is not in any file mentioned in the issue; neither
    condition navigates the inheritance/import chain deeply enough.

  \item \textbf{Semantic understanding failure} (28\% of non-sympy failures):
    the model correctly localizes the relevant file but misidentifies the
    root cause. Present in both conditions, but more common in the baseline
    (which lacks the \texttt{code\_explorer}'s iterative navigation).

  \item \textbf{Patch precision failure} (22\% of non-sympy failures):
    \texttt{str\_replace} receives an \texttt{old\_str} that does not
    match any text in the file (whitespace difference, version difference,
    or hallucinated line content). More common in the baseline condition
    (which relies solely on the pre-loaded content block).

  \item \textbf{Timeout} (15\% of non-sympy failures): the hierarchical
    system's per-instance timeout (600\,s) is exceeded on instances requiring
    large numbers of Domain Manager planning cycles.
\end{enumerate}


% ============================================================
\section{Discussion and Limitations}
\label{sec:discussion}
% ============================================================

% ============================================================
% Section 7: Discussion and Limitations
% ============================================================

\subsection{Implications for Resource-Constrained Deployment}
\label{sec:disc:deployment}

The primary practical implication of our results is that \emph{inference-time
orchestration architecture is a first-class engineering lever for improving
the capability of resource-constrained LLM deployments}. An organization
operating on a single 16\,GB GPU with a 7B model---a configuration costing
approximately \$1{,}000--\$2{,}000 in hardware---can achieve 2.21$\times$
higher code patch generation rates by adopting a hierarchical orchestration
structure rather than a monolithic single-agent loop, without any model
fine-tuning or additional hardware.

This has immediate implications for privacy-sensitive deployments (where
sending code to cloud APIs is prohibited), air-gapped environments (where
internet connectivity is absent), and cost-sensitive development workflows
(where per-token cloud API costs are prohibitive for high-volume agent use).

\subsection{Relationship to Scaling Laws}
\label{sec:disc:scaling}

Our result can be interpreted through the lens of the inference-time scaling
hypothesis: that additional compute at inference time can substitute for
additional parameters at training time~\cite{snell2024scaling,
brown2024large}. Hierarchical orchestration increases inference-time compute
in two ways: (1) by running multiple LLM calls per instance (one per agent
invocation), and (2) by generating significantly more tokens per instance
(7.7$\times$ more than the baseline). This additional compute is not
wasted---it is structural: it goes into targeted exploration (code localization),
verification (test re-execution), and plan synthesis (Domain Manager
summaries). Whether this compute investment is efficient relative to simply
using a larger model (e.g., a 70B model with equivalent VRAM via multi-GPU
setup) is an important open question.

\subsection{Limitations}
\label{sec:disc:limitations}

\paragraph{Metric.}
We report patch generation rate, not test-suite resolution rate. A system
that generates a syntactically valid non-empty patch that fails all tests
counts as a ``win'' under our metric. The gap between patch generation and
resolution is well-documented: \citet{yang2024sweagent} report that
SWE-agent generates patches on a substantially higher fraction of instances
than it resolves. Our resolution rate is thus likely lower than our $17.7\%$
patch generation figure. Future work should integrate the SWE-bench Docker
evaluation harness to report resolution rates.

\paragraph{Sympy exclusion.}
The 77 sympy instances constitute a systematic failure due to repository
setup, not model or orchestration capability. If sympy instances could be
properly evaluated, the headline numbers would shift. Our non-sympy rates
(H: 23.8\%, B: 10.8\%) may be a more representative characterization of
the system's capability on solvable instances.

\paragraph{Single model, single hardware configuration.}
All results use \texttt{Qwen2.5-Coder-7B-Instruct-AWQ} on an RTX 5080.
Generalizability to other model families (Code Llama, DeepSeek Coder,
StarCoder), quantization schemes (GPTQ, GGUF), and hardware configurations
(A100, H100, multi-GPU) is not established.

\paragraph{No ablation over hierarchy depth.}
We compare three-level (H) vs.\ one-level (B) directly, without evaluating
a two-level intermediate (orchestrator + flat workers, no domain managers).
Establishing the marginal contribution of the L2 Domain Manager layer over
a flat two-level hierarchy would require an additional experimental arm.

\paragraph{Determinism and variance.}
Each instance is run exactly once in each condition. Due to LLM sampling
non-determinism (temperature $> 0$ for most agents), there is instance-level
variance in outcomes. We do not report confidence intervals over multiple
runs due to the computational cost of re-running the full 300-instance
evaluation.

\paragraph{Wall-clock overhead.}
The hierarchical system requires an average of 94.8\,s per instance versus
5.2\,s for the baseline---an 18.2$\times$ wall-time overhead. For
interactive use cases where latency matters (e.g., IDE code completion),
this overhead is prohibitive. The hierarchical system is appropriate for
asynchronous batch workflows (e.g., nightly CI/CD bug triage, offline code
review) where throughput matters more than latency.


% ============================================================
\section{Conclusion}
\label{sec:conclusion}
% ============================================================

% ============================================================
% Section 8: Conclusion
% ============================================================

We presented \sysname, a three-level hierarchical multi-agent framework for
autonomous software engineering, and demonstrated a statistically significant
$2.21\times$ improvement in patch generation rate on SWE-bench Lite (17.7\%
vs.\ 8.0\%, McNemar $\chi^2 = 11.36$, $p = 6.0 \times 10^{-4}$) over a
controlled single-agent baseline, using an identical 7B AWQ-quantized model
on a single consumer GPU.

The core finding supports the hypothesis that the dominant failure mode for
small open-weight models on complex software engineering tasks is
\emph{context management}, not \emph{model capability}: hierarchical
decomposition into fresh, domain-scoped context windows recovers substantial
performance without any modification to model weights. The $52.8\%$
reduction in per-agent token load---despite a $7.74\times$ increase in total
pipeline token consumption---demonstrates that context isolation is a
genuine architectural improvement, not merely a reallocation of the same
work.

The Skeptical Senior Programmer verification loop---in which an independent
\texttt{test\_runner} agent re-executes the reproduction script against the
patched repository, with no access to the \texttt{patch\_writer}'s reasoning
trace---provides a mechanistic safeguard against sycophantic self-evaluation,
a failure mode documented in both the LLM reasoning
literature~\cite{huang2023self} and our own qualitative error analysis.

Future work should: (1) integrate the SWE-bench Docker harness to measure
test-suite resolution rather than patch generation; (2) resolve the sympy
repository setup failure to evaluate all 300 instances fairly; (3) evaluate
generalizability across model families and hardware configurations;
(4) add a two-level ablation arm to isolate the marginal contribution of
Domain Managers; and (5) explore training-time complementarity with
reinforcement learning approaches~\cite{luo2025swerl} that can improve
per-step reasoning quality, compounding with our inference-time architectural
gains.

The full system implementation, benchmark harness, and raw result data are
released under the GPL-3.0 license at
\texttt{https://github.com/JerMa88/subagent\_manager}.


% ============================================================
% References
% ============================================================
\bibliography{references}

\end{document}
